\begin{python0}
from solutions import *; clear()
\end{python0}
\sectionthree{Functions calling functions}

A function can call another function. For instance here's a \texttt{max3()} that uses \texttt{max2()}:

\begin{consolethree}[escapeinside=||]
#include <iostream>

int max2(int x, int y)
{
    if (x >= y)
    {
        return x;
    }
    else
    {
        return y;
    }
}

int max3(int x, int y, int z)
{
    return max2(max2(x, y), z);
}

int main()
{
    std::cout << max3(3, 5, 2) << '\n';
    return 0;
}
\end{consolethree}

In fact you have been using functions since day one: \verb!main()! is a function. When you run the program (for instance in Visual Studio .NET, you do that with Ctrl-F5), main() is the first function to execute.

Without going into details, not only can you get functions to call other functions, \texttt{main()} is called by something called the operating system. The convention is the if program Y executed successfully, the error code returned to the operating system is 0, which explains return 0 for \texttt{main()}.


